\section{Purpose and scope}\label{purpose-and-scope}

SlopNinja's primary goal is to transform writing so it sounds natural and
matches the author's intended voice. Mechanism \textit{B} serves that goal by
revising authorized text to remove unwanted formulaic prose and reduce detectable
model-writing patterns while preserving meaning, readability, and tone.
Mechanism \textit{A} supplies detection and attribution signals that support
\textit{B}'s revision search and evaluation, alongside separate preservation and
editorial checks. Together they form an adversarial training subnet:
\textit{B} trains against public \textit{A} detectors, and authorized
\textit{B} revisions challenge future \textit{A} training after evaluation
retirement. \textit{A} seeks to identify AI involvement even when \textit{B}
succeeds at reducing its visible signatures.
The subnet launches with \textit{A} alone; this development
sequence establishes a supporting evaluator and initial detection service
before the main writing product is qualified. Protocol sections state
requirements; Section~\ref{implementation-status} records which parts are
implemented today. A model, path, or component is \emph{qualified} once
recorded measurements or tests show that it meets its declared requirements;
qualification precedes activation.

\textit{A}'s origin detection estimates whether a document has human-only,
model-only, or mixed production history. Model-only and mixed work both count
as AI involvement under the declared labeling policy. The detector reports
confidence in any AI involvement and estimated AI-written or rewritten text
share as separate percentages, and attributes human and model contributors
where supported. It also locates AI-written spans of at least 24 words, with
per-sentence share and confidence (Section~\ref{span-outputs}); shorter edits
count toward both document percentages but are not localized. The scorer
serves span outputs, and once the standing cohort template enforces a span
gate, measured against SlopNinja's first span package on the same batch, every
entry must emit spans and pass it, while the reward stays document-level. Explanations
tied to particular words or constructions are later extensions.

Human writers have recurring habits of expression. Language models can also
exhibit characteristic choices in vocabulary, syntax, rhythm, and document
structure, shaped by training, tuning, and prompting. The proposed detector
learns these writing signatures through a shared origin-and-attribution model.
The hypothesis is that learning whose writing a text resembles helps determine
whether AI was involved. Signatures can overlap and change with prompts or
editing; they are statistical evidence, not unique identifiers. Their value
for AI detection must be measured on unseen writers, generators, and workflows.

\textit{A} learns these signatures from raw text and uses no grammar features.
A shared English grammar is planned as a later input for \textit{B}, where it
locates typed word, phrase, and clause edits and guides revision. Detector
evaluation can begin before transformation, private serving, or grammar
incentives are ready.

The three logical mechanisms have distinct jobs:
\begin{itemize}\tightlist
\item \textit{A}, AI detection and attribution: one text model estimates AI
involvement from production origin, matches human or AI-model profiles, and
locates AI spans of at least 24 words. Entrant models are
published, licensed CC BY-SA 4.0, when each daily evaluation cohort's scoring
starts.
Launch includes \textit{A} alone.
\item \textit{B}, authorized transformation: a private model proposes revisions
that improve fit to a writer's brief and reduce detector scores while preserving
meaning, readability, and intended tone. It activates separately.
\item \textit{C}, grammar improvement for \textit{B}: accepted reusable lexical
and structural contributions to the grammar that \textit{B} uses for edits earn
credit once, after a separate acceptance and reward protocol is qualified. It
activates by its own decision, possibly before \textit{B}.
\end{itemize}
These names describe workstreams, not three available native chain mechanisms.
The launch uses one native mechanism; later allocation must fit the chain's
supported mechanism interface (Section~\ref{benchmark-scoring-and-mechanism-mapping}).

\subsection{What a prediction means}

Origin is a production-history claim. Human-only, model-only, and mixed labels
refer to how a document was made under an explicit labeling policy. Attribution
instead asks which registered human or model profiles it resembles. A known
writer may use AI assistance; an unfamiliar model may produce text with no good
gallery match. Neither a matching human profile nor a low detector score proves
human-only origin. An unavailable prediction must remain an error or abstention,
never a default human label.

The launch evaluates English documents within declared length and resource
bounds. The evidence so far comes from passages whose human documents are
whole-paragraph excerpts of 150 to 450 words; other lengths need their own
measurements. Span outputs rely on retained exact span labels; multilingual
detection and manuscript-scale support each require suitable data and
measurements. Handling an embedded
foreign phrase, equation, or code block does not establish detection quality
for that language or content type. Predictions must retain uncertainty,
calibration context, and the limitations of the measured population.

We use \emph{slop} for defects in prose: vague or formulaic phrasing, needless
repetition, filler, and other patterns that weaken the piece. This is an
editorial definition, separate from production origin. Informal language,
unusual grammar, or a writer's recurring verbal habit may be appropriate to a
brief. The shared grammar planned for \textit{B}'s edits likewise describes
English as used, including errors and nonstandard forms, rather than enforcing
a prescriptive ideal.

\subsection{The grammar's role in later editing}\label{grammar-role}

A resolved parse exposes lexical choices, grammatical relations, and document
structure in a reusable representation. \textit{B} is planned to use the shared
English grammar to divide a source into typed word, phrase, and clause edit
units and to compare each revision with its source (Section~\ref{later-b}).
\textit{A} does not use it. The first real-data test of a CFG-only detector
failed, and adding CFG features to text gave only a small, uncertain gain, so
\textit{A}'s input contract is raw text (Section~\ref{input-contract}). A
head-to-head test of writer profiles found CFG writer vectors below
\textit{A}'s encoder on writer verification, by \(-0.034\) in AUC, simple
stylometric features above both, and only \(+0.0015\) from adding CFG to the
encoder plus stylometry, so private writer profiles use the encoder and
stylometric features (Section~\ref{writer-profile-evidence}). For \textit{B},
parsing costs money and time, and resolution can be wrong. Whether grammar-typed edits help
\textit{B}, whether \textit{C}'s contributions improve them, and any later
explanations each need their own evidence.

\subsection{Commercial thesis and users}\label{commercial-thesis}

The first paid product is hosted AI detection for applications, publishers,
researchers, and document workflows that need to assess AI involvement in
writing, with origin probabilities and attribution as additional detail.
\textit{A} artifacts become freely executable when their cohort's scoring starts. Paid
access, through the detection API or quoted jobs, offers maintained attribution profiles, bounded execution of a qualified
artifact, private delivery, and integration convenience. An \textit{A} detection job
runs the raw-text model alone, with no grammar parse. Publishing a model does not establish
that users will pay for hosting it; the commercial pilot must measure repeat
use, willingness to pay, false-positive disputes, and delivery cost.

Testnet runs for several months before the mainnet subnet is registered.
During that period SlopNinja plans to earn revenue from free,
advertising-supported public detection pages served by \textit{A} and from
other products, accruing in a publicly visible pot (Section~\ref{testnet-revenue}).
The pot pays the mainnet subnet's registration once it holds about 1.5 times
its cost. After the launch, the rest buys the subnet's alpha: half is burned,
and half is distributed to testnet miners in proportion to the revenue earned
by their winning models, those that beat the reference; all later revenue buys
and burns alpha.
The detection page itself carries no advertising and stores no text; ads run
on the site's content pages, and a paid detection API follows once the served package is
published. The pot is a published multisignature TAO account with a public
ledger of every payout, conversion, purchase, burn, and claim.
The coordinator decides when measured quality is sufficient to open sales and
which workloads may be sold. Benchmark eligibility alone does not enable
checkout. The initial price hypothesis aims below Pangram's standalone API
rate, subject to an independently evaluated quality comparison and measured
costs. The principal of a quoted paid job buys and burns alpha, while the
detection API's revenue joins the pot. Spendable operating
funds and refund reserves must come from an explicitly funded treasury,
independently of the burn (Section~\ref{revenue-and-costs}).

The main intended writing product, activated after the initial detection service,
serves users who supply samples of their own writing. A private writer profile
summarizes those samples from \textit{A}'s encoder and stylometric features;
\textit{A} compares successive revisions with that profile, and \textit{B} uses
its author-fit signal to guide editing toward the requested voice. No profile
tested so far ranks rewrites toward a writer reliably, so this workflow waits
for a learned matcher qualified on fresh writers
(Section~\ref{author-and-origin-detection}). It aims to
preserve every argument, qualification, attribution, and necessary detail while
improving voice and removing unwanted formulaic prose. Editorial quality and
production history remain separate: a revision of model-generated writing
retains that history even if readers prefer it or detectors miss it. Incremental
improvement and attainment of a strict detector target are also separate claims
(Section~\ref{transformation-and-improvement-utility}).

The intended adversarial training loop uses public detectors during
private revision search and reuses authorized revisions after evaluation
retirement. It must measure whether adaptive editing improves prose or merely
exploits the current judges. Writer consent and source rights limit training
reuse; customer jobs do not become benchmark material by default. The same
tools can be used to obscure provenance or impersonate a writer. Product claims
must preserve that risk and must not represent detector evasion as evidence of
human authorship.
